\medskip
\noindent\textbf{Case (3): Artin--Schreier--Witt.}
A natural alternative to the direct calculation below would be to filter a
cyclic $p$-power torsor by its degree-$p$ layers and try to apply the
Artin--Schreier calculation inductively. In
Appendix~\ref{appendix:artin_schreier_induction_failure} we examine this approach and
explain why the two expected induction inputs do not by themselves control
ramification before the intervening wild base change.

If the connected inertia group $H$ produced by the common reduction has order
$p$, the Artin--Schreier calculation in Case~(2) already gives the required
unramified conclusion. We may therefore assume that $H$ has order $p^r$ with
$r\geq2$, and continue to denote it by $G$.
By \Cref{lemma:realization_over_Gm}(3), the connected torsor has an equation
\[
    \wp\vec X=\vec f=(f_0,\dots,f_{r-1}),
    \qquad f_j\in x^{-1}k[x^{-1}],
\]
where, putting $m_j=\deg_{x^{-1}}f_j$ for $f_j\neq0$ and $m_j=0$ for
$f_j=0$, the pole bounds are
\begin{equation}\label{eq:witt_pole_bound_in_calculation}
    p^{r-1-j}m_j<d\qquad(0\leq j<r).
\end{equation}
Put $R=k[x,x^{-1}]$ and
\[
    S=R[\vec X]/(\wp\vec X\ominus\vec f).
\]
After pulling back to $\bG_m^d$, define the single Witt vectors
\[
    \vec Y:=\bigoplus_{i=1}^d\vec X_i,
    \qquad
    \vec\alpha:=\bigoplus_{i=1}^d\vec f(x_i).
\]
Here and below $\bigoplus$ denotes iterated Witt addition, not a direct sum.

\begin{prop}[Identification of the Witt product torsor]
\label{prop:witt_cover_identification}
Every component of $\vec Y$ is $\KG$-invariant, every component of
$\vec\alpha$ is a symmetric Laurent polynomial, and
\[
    (S^{\otimes_k d})^{\KG}
       \cong
    R^{\otimes_k d}[\vec Y]/(\wp\vec Y\ominus\vec\alpha)
\]
as $G$-torsors over $R^{\otimes_k d}$. Taking $S_d$-invariants gives
\[
    \bigl((S^{\otimes_k d})^{\KG}\bigr)^{S_d}
       \cong
    k[e_1,\dots,e_d,e_d^{-1}][\vec Y]
       /(\wp\vec Y\ominus\vec\alpha).
\]
Consequently,
\[
    \cP^{[d]}\big|_{\Spec\widehat K_\eta}
       \cong
    \Spec\widehat K_\eta[\vec Y]
       /(\wp\vec Y\ominus\vec\alpha).
\]
\end{prop}

\begin{proof}
Write an element of $\KG\cong G^{d-1}$ as
$(g_1,\dots,g_{d-1})$. Under the usual embedding in $G^d$, it acts in Witt
coordinates by
\[
\begin{aligned}
    \vec X_1&\longmapsto\vec X_1\oplus\underline g_1,\\
    \vec X_i&\longmapsto
       \vec X_i\oplus\underline g_i\ominus\underline g_{i-1}
       &&(1<i<d),\\
    \vec X_d&\longmapsto\vec X_d\ominus\underline g_{d-1},
\end{aligned}
\]
where $\underline g$ is the image of $g$ in $W_r(\F_p)$. Functoriality of
$W_r$ means that this action is componentwise and commutes with $\oplus$ and
$\ominus$. The translations telescope in $\vec Y$, so each component of
$\vec Y$ is invariant. Additivity of $\wp$ gives
\[
    \wp\vec Y
       =\bigoplus_{i=1}^d\wp\vec X_i
       =\bigoplus_{i=1}^d\vec f(x_i)
       =\vec\alpha.
\]
The resulting coordinate map is equivariant for the residual $G$-action.
Both the displayed Artin--Schreier--Witt cover and the contracted product are
$G$-torsors over $R^{\otimes_k d}$, by
\Cref{lemma:asw_etale_torsor,prop:parallel_cyclic_product_torsor}; hence the
map is an isomorphism by \Cref{prop:torsor_morphism_iso}.

The group $S_d$ permutes the tensor factors. Commutativity and associativity
of Witt addition show that it fixes $\vec Y$ and $\vec\alpha$ componentwise.
Thus every component of $\vec\alpha$ lies in
\[
    (R^{\otimes_k d})^{S_d}=k[e_1,\dots,e_d,e_d^{-1}].
\]
By \Cref{lemma:asw_etale_torsor}, the ASW algebra is free over
$R^{\otimes_k d}$ with basis
$\prod_nY_n^{a_n}$, $0\leq a_n<p$. Every basis element is $S_d$-invariant,
so an element is invariant exactly when all its coefficients are invariant.
This proves the second displayed isomorphism. Finally,
\Cref{prop:symmetric_power_semidirect_quotient} identifies the quotient by
$\KG\rtimes S_d$ with $\cP^{[d]}$, and base change to $\widehat K_\eta$
gives the last assertion.
\end{proof}

It remains to show that every component of $\vec\alpha$ is regular at the
exceptional valuation. The Witt components are not ordinary componentwise
sums, so this is the weighted analogue of the Artin--Schreier degree count.
We claim that every component $\alpha_n$, $0\leq n<r$, belongs to $\kappa$.

Give $f_j(x_i)$ weight $p^j$. By
\Cref{lemma:witt_addition_isobaric}, every nonzero monomial
\[
    \prod_{i,j}f_j(x_i)^{a_{ij}}
\]
appearing in $\alpha_n$ satisfies
\[
    \sum_{i,j}a_{ij}p^j=p^n.
\]
Put $A_j:=\sum_i a_{ij}$. Since $f_j(x_i)$ is a polynomial in
$y_i=x_i^{-1}$ of degree $m_j$, every monomial obtained after expansion in
the $y_i$ has total degree at most
\[
    \sum_jA_jm_j
       <d\sum_jA_jp^{j-(r-1)}
       =d\,p^{n-(r-1)}
       \leq d.
\]
The inequality is strict: a nonzero monomial has a positive exponent on at
least one nonzero $f_j$, and every pole bound in
\Cref{eq:witt_pole_bound_in_calculation} is strict. By
\Cref{prop:witt_cover_identification}, the whole component $\alpha_n$ is
symmetric. Therefore \Cref{lemma:symmetric_regularity} gives
$\alpha_n\in\kappa$.
Thus
\[
    \vec\alpha\in W_r(\kappa)\subset W_r(\widehat\Oo_\eta).
\]

It follows that the equation
\[
    \wp\vec Y=\vec\alpha
\]
is defined over $\widehat\Oo_\eta$. By
\Cref{lemma:asw_etale_torsor}, it defines a finite \'etale
$\widehat\Oo_\eta$-algebra. Its generic fiber is the restriction of
$\cP^{[d]}$ to $\widehat K_\eta$ by
\Cref{prop:witt_cover_identification}. Hence every field factor is
unramified. This proves part~(3).
